% This file is generated from the corresponding Obsidian note.
% Source: Teoricos/Glosario Teorico.md
\chapter{Guía de lectura}
\label{ch:glosario-teorico}
\section{Indice de teoricos}
\subsection{Teoricos 15-25}
\begin{itemize}
\item \hyperref[ch:teo-15]{Teórico 15}: campos vectoriales suaves, campos coordenados, extension de vectores tangentes, campos en subvariedades, curvas integrales y teorema fundamental de EDO autonomas.
\item \hyperref[ch:teo-16]{Teórico 16}: existencia y unicidad de curvas integrales, flujo maximal, flujo global, campos completos y mapas de tiempo.
\item \hyperref[ch:teo-17]{Teórico 17}: propiedades del flujo, generador infinitesimal, uniformidad global del tiempo, campos con soporte compacto y lema del escape.
\item \hyperref[ch:teo-18]{Teórico 18}: variedades paralelizables, extension local de campos, aplicaciones del flujo y teorema de enderezamiento.
\item \hyperref[ch:teo-19]{Teórico 19}: subgrupos monoparametricos, corchete de Lie, campos relacionados, derivada de Lie, conmutacion de flujos y rectificacion simultanea.
\item \hyperref[ch:teo-20]{Teórico 20}: algebra multilineal, tensores covariantes, tensores alternantes, alternador, producto exterior y algebra exterior.
\item \hyperref[ch:teo-21]{Teórico 21}: fibrado de formas alternantes, suavidad de formas, producto cuña, algebra de formas y diferencial exterior.
\item \hyperref[ch:teo-22]{Teórico 22}: complejo de de Rham, formas cerradas y exactas, pullback de formas.
\item \hyperref[ch:teo-23]{Teórico 23}: distribuciones suaves, distribuciones integrables, involutividad y teorema de Frobenius local.
\item \hyperref[ch:teo-24]{Teórico 24}: orientacion de espacios vectoriales, orientacion de variedades, cartas orientadas, orientabilidad de hipersuperficies y orientacion inducida en el borde.
\item \hyperref[ch:teo-25]{Teórico 25}: medida nula, dominios de integracion, integrales de formas en abiertos y variedades, orientacion de formas, difeomorfismos que preservan orientacion e integracion sobre parametrizaciones.
\end{itemize}
\section{Teorico 18}
\begin{booklemma}{Otra forma de hacerlo (creo)}
Sea $M$ una variedad suave de dimensión $n$ y sea $p\in M$. Si $\{v_1,\ldots,v_k\}\subset T_pM$ es linealmente independiente, entonces existe una carta suave $(U,\varphi=(x_1,\ldots,x_n))$ centrada en $p$ tal que

\[
\left.\frac{\partial}{\partial x_j}\right|_p=v_j,\qquad j=1,\ldots,k.
\]
\end{booklemma}
\begin{bookproof}{}
\begin{enumerate}
\item Completamos $\{v_1,\ldots,v_k\}$ a una base $\{v_1,\ldots,v_n\}$ de $T_pM$.
\item Tomemos una carta suave $(U,\psi=(y_1,\ldots,y_n))$ centrada en $p$, es decir, $\psi(p)=0$.
\item Como $(d\psi)_p:T_pM\to T_0\mathbb R^n$ es un isomorfismo, los vectores
\end{enumerate}
\[
w_j:=(d\psi)_p(v_j),\qquad j=1,\ldots,n,
\]
forman una base de $T_0\mathbb R^n$. 4. Identificamos canónicamente $T_0\mathbb R^n$ con $\mathbb R^n$ mediante

\[
\sum_{i=1}^n a_i\left.\frac{\partial}{\partial r_i}\right|_0\longleftrightarrow (a_1,\ldots,a_n).
\]
Denotemos por $\widetilde w_j\in\mathbb R^n$ al vector correspondiente a $w_j$ bajo esta identificación. 5. Como $\{w_1,\ldots,w_n\}$ es base de $T_0\mathbb R^n$, entonces $\{\widetilde w_1,\ldots,\widetilde w_n\}$ es base de $\mathbb R^n$. 6. Definimos $T:\mathbb R^n\to\mathbb R^n$ como la única transformación lineal tal que

\[
T(\widetilde w_j)=e_j,\qquad j=1,\ldots,n,
\]
donde $\{e_1,\ldots,e_n\}$ es la base estándar de $\mathbb R^n$. 7. Como $T$ manda una base en una base, $T$ es invertible. Por lo tanto $T$ es un difeomorfismo lineal de $\mathbb R^n$ en $\mathbb R^n$. 8. Definimos ahora

\[
\varphi=T\circ\psi.
\]
Como $\psi$ es carta suave y $T$ es un difeomorfismo lineal, $\varphi$ es una carta suave. 9. Además, como $\psi(p)=0$ y $T$ es lineal, tenemos

\[
\varphi(p)=T(\psi(p))=T(0)=0.
\]
Luego $\varphi$ está centrada en $p$. 10. Veamos qué hace su diferencial sobre los $v_j$. Por regla de la cadena,

\[
(d\varphi)_p(v_j)=(dT)_{\psi(p)}((d\psi)_p(v_j))=(dT)_0(w_j).
\]
\begin{enumerate}[start=11]
\item Como $T$ es lineal, $(dT)_0$ actúa sobre los vectores tangentes usando la misma transformación lineal $T$ sobre sus coordenadas. Como $T(\widetilde w_j)=e_j$, se obtiene
\end{enumerate}
\[
(dT)_0(w_j)=\left.\frac{\partial}{\partial r_j}\right|_0.
\]
\begin{enumerate}[start=12]
\item Por lo tanto,
\end{enumerate}
\[
(d\varphi)_p(v_j)=\left.\frac{\partial}{\partial r_j}\right|_0.
\]
\begin{enumerate}[start=13]
\item Por otro lado, si $\varphi=(x_1,\ldots,x_n)$, entonces por definición del vector coordenado asociado a la carta $\varphi$,
\end{enumerate}
\[
(d\varphi)_p\left(\left.\frac{\partial}{\partial x_j}\right|_p\right)=\left.\frac{\partial}{\partial r_j}\right|_{\varphi(p)}.
\]
\begin{enumerate}[start=14]
\item Como $\varphi(p)=0$, esto queda
\end{enumerate}
\[
(d\varphi)_p\left(\left.\frac{\partial}{\partial x_j}\right|_p\right)=\left.\frac{\partial}{\partial r_j}\right|_0.
\]
\begin{enumerate}[start=15]
\item Comparando con el paso 12, obtenemos
\end{enumerate}
\[
(d\varphi)_p(v_j)=(d\varphi)_p\left(\left.\frac{\partial}{\partial x_j}\right|_p\right).
\]
\begin{enumerate}[start=16]
\item Como $(d\varphi)_p:T_pM\to T_0\mathbb R^n$ es un isomorfismo, concluimos que
\end{enumerate}
\[
v_j=\left.\frac{\partial}{\partial x_j}\right|_p,\qquad j=1,\ldots,n.
\]
\begin{enumerate}[start=17]
\item En particular, esto vale para $j=1,\ldots,k$, que era lo que queríamos probar.
\end{enumerate}
\bookqed
\end{bookproof}
\phantomsection\label{blk:3071e3}
\begin{bookproposition}{}
Sea $\sigma:U\subseteq\mathbb R^k\to M$ una aplicación suave, con $0\in U$, y sea $\alpha:\mathbb R\to U$ la curva dada por $\alpha(s)=se_1=(s,0,\ldots,0)$. Entonces

\[
(\sigma\circ\alpha)'(0)=(d\sigma)_0\left(\frac{\partial}{\partial r^1}\big|_0\right).
\]
\end{bookproposition}
\begin{bookproof}{}
\begin{enumerate}
\item Por definición de velocidad de una curva, tenemos
\end{enumerate}
\[
\alpha'(0)=(d\alpha)_0\left(\frac{\partial}{\partial s}\big|_0\right)\in T_0\mathbb R^k
\]
\begin{enumerate}[start=2]
\item Queremos probar que
\end{enumerate}
\[
\alpha'(0)=\frac{\partial}{\partial r^1}\big|_0
\]
\begin{enumerate}[start=3]
\item Para eso basta probar que ambos vectores tangentes actúan igual sobre toda función $g\in C^\infty(\mathbb R^k)$.
\item Por definición de diferencial,
\end{enumerate}
\[
\alpha'(0)(g)=(d\alpha)_0\left(\frac{\partial}{\partial s}\big|_0\right)(g)=\frac{\partial}{\partial s}\big|_0(g\circ\alpha)
\]
\begin{enumerate}[start=5]
\item Como $\alpha(s)=(s,0,\ldots,0)$, se tiene
\end{enumerate}
\[
\frac{\partial}{\partial s}\big|_0(g\circ\alpha)=\frac{d}{ds}\big|_{s=0}g(s,0,\ldots,0)
\]
\begin{enumerate}[start=6]
\item Pero esto es exactamente la derivada parcial de $g$ respecto de la primera coordenada en $0$, es decir,
\end{enumerate}
\[
\frac{d}{ds}\big|_{s=0}g(s,0,\ldots,0)=\frac{\partial g}{\partial r^1}(0)
\]
\begin{enumerate}[start=7]
\item Por definición del vector coordenado,
\end{enumerate}
\[
\frac{\partial}{\partial r^1}\bigg|_0(g)=\frac{\partial g}{\partial r^1}(0)
\]
\begin{enumerate}[start=8]
\item Luego, para toda $g\in C^\infty(\mathbb R^k)$, se cumple
\end{enumerate}
\[
\alpha'(0)(g)=\frac{\partial}{\partial r^1}\big|_0(g)
\]
\begin{enumerate}[start=9]
\item Por lo tanto,
\end{enumerate}
\[
\alpha'(0)=\frac{\partial}{\partial r^1}\big|_0
\]
\begin{enumerate}[start=10]
\item En consecuencia, si $\sigma:U\subseteq\mathbb R^k\to M$ es suave y $0\in U$, entonces
\end{enumerate}
\[
(\sigma\circ\alpha)'(0)=d(\sigma\circ\alpha)_0\left(\frac{\partial}{\partial s}\big|_0\right)=(d\sigma)_0\left((d\alpha)_0\left(\frac{\partial}{\partial s}\big|_0\right)\right)=(d\sigma)_0\left(\frac{\partial}{\partial r^1}\big|_0\right)
\]
\bookqed
\end{bookproof}
\phantomsection\label{blk:99bdaa}
\section{Teo 21}
\begin{bookremark}{}
Sea $f\in C^{\infty}(M)$ y $v\in T_{p}M$ vector tangente entonces

\[
(df)_{p}(v)=v(f)
\]
\end{bookremark}
\begin{bookproof}{}
\begin{enumerate}
\item Como $f:M\rightarrow\mathbb{R}$ entonces $(df)_{p}:T_{p}M\rightarrow T_{f(p)}\mathbb{R}$
\item Luego
\end{enumerate}
\[
(df)_{p}(v)=c(p)\frac{\partial}{\partial t}|_{f(p)}
\]
\begin{enumerate}[start=3]
\item Entonces $c(p)=(df)_{p}(v)(t)=v(t\circ f)=v(f)$
\item Luego
\end{enumerate}
\[
(df)_{p}(v)=v(f)\frac{\partial}{\partial t}|_{f(p)}\simeq v(f)
\]
donde use la identificacion $T_{f(p)}\mathbb{R}\simeq \mathbb{R}$ 5. Por lo tanto $df(v)=v(f)$

\bookqed
\end{bookproof}
\begin{bookremark}{}
Sea $f\in C^{\infty}(M)$ y sea $v\in T_pM$. Entonces, bajo la identificación $T_{f(p)}\mathbb R\simeq \mathbb R$, se tiene

\[
(df)_p(v)=v(f).
\]
\end{bookremark}
\begin{bookproof}{}
\begin{enumerate}
\item Como $f:M\rightarrow\mathbb{R}$, entonces
\end{enumerate}
\[
(df)_p:T_pM\rightarrow T_{f(p)}\mathbb{R}.
\]
\begin{enumerate}[start=2]
\item Luego existe $c\in\mathbb R$ tal que
\end{enumerate}
\[
(df)_p(v)=c\frac{\partial}{\partial t}\bigg|_{f(p)}.
\]
\begin{enumerate}[start=3]
\item Entonces
\end{enumerate}
\[
c=(df)_p(v)(t)=v(t\circ f)=v(f),
\]
porque $t:\mathbb R\to\mathbb R$ es la coordenada identidad. 4. Luego

\[
(df)_p(v)=v(f)\frac{\partial}{\partial t}\bigg|_{f(p)}\simeq v(f),
\]
donde usamos la identificación $T_{f(p)}\mathbb{R}\simeq \mathbb{R}$ dada por

\[
a\frac{\partial}{\partial t}\bigg|_{f(p)}\mapsto a.
\]
\begin{enumerate}[start=5]
\item Por lo tanto, bajo esta identificación,
\end{enumerate}
\[
(df)_p(v)=v(f).
\]
\bookqed
\end{bookproof}
\phantomsection\label{blk:311f86}
\begin{bookproposition}{Hulet 9.15}
Sean $X,Y\in\mathfrak{X}(M)$ tales que $[X,Y]=0$, y sean $\theta,\phi$ los flujos de $X$ e $Y$, respectivamente. Dado $p\in M$, sean $U,\delta$ dados por el lema de uniformidad conjunta. Entonces:

\begin{enumerate}
\item $d\theta_t\circ Y=Y\circ\theta_t$, en $U$, para todo $|t|<\delta$.
\item $\theta_t\circ\phi_s=\phi_s\circ\theta_t$, en $U$, para todos $s,t$ con $|s|,|t|<\delta$.
\end{enumerate}
\end{bookproposition}
\begin{bookproof}{}
Recordemos la fórmula

\[
[X,Y](p)=\frac{d}{dt}\bigg|_{t=0}(d\theta_{-t})_{\theta_{t}(p)}\left(Y_{\theta_t(p)}\right),\qquad p\in M.
\]
\begin{itemize}
\item (1)
\end{itemize}
\begin{enumerate}
\item Probar $d\theta_t\circ Y=Y\circ\theta_t$ equivale a probar que
\end{enumerate}
\[
Y_p=d\theta_{-t}\left(Y_{\theta_t(p)}\right),\qquad \forall t,\forall p.
\]
\begin{enumerate}[start=2]
\item En efecto, la igualdad
\end{enumerate}
\[
d\theta_t\circ Y=Y\circ\theta_t
\]
significa que para todo $p$ se tiene

\[
(d\theta_t)_p(Y_p)=Y_{\theta_t(p)}.
\]
Como $\theta_{-t}$ es la inversa de $\theta_t$, aplicando $(d\theta_{-t})_{\theta_t(p)}$ a ambos lados obtenemos

\[
Y_p=(d\theta_{-t})_{\theta_t(p)}(Y_{\theta_t(p)}).
\]
(Notar que $(d\theta_{t})_{p}:T_{p}M\rightarrow T_{\theta_{t}(p)}M$ por eso tomo $(d\theta_{-t})_{\theta_t(p)}$ para que quede bien definida la composicion de la izquierda y obviaemnte $\theta_{-t}\circ\theta_{t}=Id$ y el diferencial de la identidad sigue siendo la identidad) 3. Recíprocamente, si vale esta última igualdad, aplicando $(d\theta_t)_p$ a ambos lados se recupera

\[
(d\theta_t)_p(Y_p)=Y_{\theta_t(p)}.
\]
Por lo tanto ambas igualdades son equivalentes. 4. En otras palabras, para probar la primera parte, basta ver que para cada $p$, la curva

\[
z(t)=(d\theta_{-t})_{\theta_{t}(p)}\left(Y_{\theta_t(p)}\right)
\]
es constantemente igual a $Y_p$. 5. Calculamos $z'(s)$. Primero, usando el cambio $t\mapsto t+s$,

\[
z'(s)=\frac{d}{dt}\bigg|_{t=0}z(t+s)=\frac{d}{dt}\bigg|_{t=0}(d\theta_{-t-s})_{\theta_{t+s}(p)}\left(Y_{\theta_{t+s}(p)}\right).
\]
\begin{enumerate}[start=6]
\item Como $\theta_{t+s}(p)=\theta_t(\theta_s(p))$ y $\theta_{-t-s}=\theta_{-s}\circ\theta_{-t}$, tenemos
\end{enumerate}
\[
z'(s)=\frac{d}{dt}\bigg|_{t=0}(d\theta_{-s})_{\theta_s(p)}\left((d\theta_{-t})_{\theta_{t}(\theta_{s}(p))}\left(Y_{\theta_t(\theta_s(p))}\right)\right).
\]
\begin{enumerate}[start=7]
\item Como $(d\theta_{-s})_{\theta_s(p)}$ no depende de $t$, podemos sacarlo fuera de la derivada:
\end{enumerate}
\[
z'(s)=(d\theta_{-s})_{\theta_s(p)}\left(\frac{d}{dt}\bigg|_{t=0}(d\theta_{-t})_{\theta_{t}(\theta_{s}(p))}\left(Y_{\theta_t(\theta_s(p))}\right)\right).
\]
\begin{enumerate}[start=8]
\item Por la fórmula del corchete aplicada en el punto $\theta_s(p)$, se obtiene
\end{enumerate}
\[
z'(s)=(d\theta_{-s})_{\theta_s(p)}\left([X,Y]_{\theta_s(p)}\right).
\]
\begin{enumerate}[start=9]
\item Como $[X,Y]=0$, resulta
\end{enumerate}
\[
z'(s)=0,\qquad \forall s.
\]
\begin{enumerate}[start=10]
\item Luego $z$ es constante. Y considerando $\theta_{0}=Id_{M}$ entonces $(d\theta_{0})_{p}=Id_{T_{p}M}$ tenemos
\end{enumerate}
\[
z(0)=(d\theta_0)_{\theta_{0}(p)}(Y_{\theta_0(p)})=Y_p,
\]
concluimos que

\[
z(s)=Y_p,\qquad \forall s.
\]
\begin{enumerate}[start=11]
\item Por lo tanto
\end{enumerate}
\[
d\theta_{-s}\left(Y_{\theta_s(p)}\right)=Y_p\quad\forall t,\forall p
\]
como queriamos

\begin{itemize}
\item (2)
\end{itemize}
\begin{enumerate}
\item Ahora probamos la segunda parte. Fijado $p\in U$, debemos probar que
\end{enumerate}
\[
\phi_s(p)=\theta_{-t}\circ\phi_s\circ\theta_t(p),\qquad \forall s.
\]
\begin{enumerate}[start=2]
\item Es suficiente ver que el lado derecho es una curva integral de $Y$ que para $s=0$ pasa por $p$.
\item Definimos
\end{enumerate}
\[
\sigma(s)=\theta_{-t}\circ\phi_s\circ\theta_t(p).
\]
\begin{enumerate}[start=4]
\item Entonces
\end{enumerate}
\[
\sigma(0)=\theta_{-t}\circ\phi_0\circ\theta_t(p)=\theta_{-t}\circ\theta_t(p)=p.
\]
\begin{enumerate}[start=5]
\item Además, aplicando la regla de la cadena tenemos
\end{enumerate}
\[
\sigma'(s)=(d\theta_{-t})_{\phi_s(\theta_t(p))}\left(\frac{d}{du}\bigg|_{u=s}\phi_u(\theta_t(p))\right).
\]
\begin{enumerate}[start=6]
\item Como $\phi$ es el flujo de $Y$, para cada punto fijo $q\in M$ la curva
\end{enumerate}
\[
\gamma(u)=\phi_u(q)
\]
es curva integral de $Y$. Tomando

\[
q=\theta_t(p),
\]
tenemos

\[
\gamma(u)=\phi_u(\theta_t(p)).
\]
Entonces

\[
\gamma'(s):= (d\gamma)_s\left(\frac{d}{du}\bigg|_s\right)=Y_{\gamma(s)}=Y_{\phi_s(\theta_t(p))}.
\]
Es decir,

\[
\frac{d}{du}\bigg|_{u=s}\phi_u(\theta_t(p)):=\gamma'(s)=Y_{\phi_s(\theta_t(p))}.
\]
\begin{enumerate}[start=7]
\item Por lo tanto
\end{enumerate}
\[
\sigma'(s)=(d\theta_{-t})_{\phi_s(\theta_t(p))}\left(Y_{\phi_s(\theta_t(p))}\right).
\]
\begin{enumerate}[start=8]
\item Usando la primera parte con tiempo $-t$, obtenemos
\end{enumerate}
\[
\sigma'(s)=Y_{\theta_{-t}(\phi_s(\theta_t(p)))}=Y_{\sigma(s)}.
\]
\begin{enumerate}[start=9]
\item Así, $\sigma$ es una curva integral de $Y$ que empieza en $p$.
\item Por unicidad de curvas integrales,
\end{enumerate}
\[
\sigma(s)=\phi_s(p).
\]
\begin{enumerate}[start=11]
\item Luego
\end{enumerate}
\[
\phi_s(p)=\theta_{-t}\circ\phi_s\circ\theta_t(p).
\]
\begin{enumerate}[start=12]
\item Componiendo a la izquierda con $\theta_t$, se obtiene
\end{enumerate}
\[
\theta_t\circ\phi_s(p)=\phi_s\circ\theta_t(p).
\]
\begin{enumerate}[start=13]
\item Como esto vale para todo $p\in U$, concluimos que
\end{enumerate}
\[
\theta_t\circ\phi_s=\phi_s\circ\theta_t
\]
en $U$, para todos $s,t$ con $|s|,|t|<\delta$.

\bookqed
\end{bookproof}
\begin{bookproposition}{Hulet 9.16}
Sean $X,Y\in\mathfrak{X}(M)$, cuyos flujos conmutan. Entonces

\[
[X,Y]=0.
\]
\end{bookproposition}
\begin{bookproof}{}
\begin{enumerate}
\item Sean $\theta$ y $\phi$ los flujos de $X$ e $Y$, respectivamente. Fijemos $p\in M$. Por el lema de uniformidad conjunta de los flujos, existen un abierto $U$ que contiene a $p$ y $\delta>0$ tales que, para $|s|,|t|<\delta$, las composiciones $\theta_t\circ\phi_s$ y $\phi_s\circ\theta_t$ están definidas en $U$.
\item Como los flujos conmutan, para $|s|,|t|<\delta$ se cumple
\end{enumerate}
\[
\theta_t(\phi_s(p))=\phi_s(\theta_t(p)).
\]
\begin{enumerate}[start=3]
\item Derivando respecto de $s$ en $s=0$, obtenemos
\end{enumerate}
\[
d\theta_{t}(Y_{p})=d\theta_t\left(\frac{d}{ds}\bigg|_{s=0}\phi_s(p)\right)=\frac{d}{ds}\bigg|_{s=0}\phi_s(\theta_t(p))=Y_{\theta_{t}(p)}
\]
\begin{enumerate}[start=4]
\item Aplicamos $(d\theta_{-t})_{\theta_t(p)}$ a ambos lados. Entonces
\end{enumerate}
\[
\begin{aligned}(d\theta_{-t})_{\theta_t(p)}\left((d\theta_t)_p(Y_p)\right)&=(d\theta_{-t})_{\theta_t(p)}\left(Y_{\theta_t(p)}\right),\\d(\theta_{-t}\circ\theta_t)_p(Y_p)&=(d\theta_{-t})_{\theta_t(p)}\left(Y_{\theta_t(p)}\right),\\Y_p&=(d\theta_{-t})_{\theta_t(p)}\left(Y_{\theta_t(p)}\right).\end{aligned}
\]
\begin{enumerate}[start=5]
\item Por lo tanto, la curva
\end{enumerate}
\[
t\mapsto(d\theta_{-t})_{\theta_t(p)}\left(Y_{\theta_t(p)}\right)
\]
es constante para $|t|<\delta$. 6. Usando la fórmula del corchete en términos del flujo de $X$,

\[
[X,Y]_p=\frac{d}{dt}\bigg|_{t=0}(d\theta_{-t})_{\theta_t(p)}\left(Y_{\theta_t(p)}\right)=0.
\]
\begin{enumerate}[start=7]
\item Como $p\in M$ era arbitrario, concluimos que
\end{enumerate}
\[
[X,Y]=0.
\]
\bookqed
\end{bookproof}
\begin{booktheorem}{}
Sean $X_1,\ldots,X_k\in\mathfrak X(M)$ campos linealmente independientes en un entorno de $p\in M$, y supongamos que

\[
[X_i,X_j]=0,\qquad 1\leq i,j\leq k.
\]
Entonces existe una carta

\[
(U,\varphi=(x^1,\ldots,x^n))
\]
alrededor de $p$ tal que

\[
X_i|_U=\frac{\partial}{\partial x^i},\qquad i=1,\ldots,k.
\]
\end{booktheorem}
\begin{bookproof}{}
\begin{enumerate}
\item Sea $n=\dim M$. Como
\end{enumerate}
\[
X_1(p),\ldots,X_k(p)
\]
son linealmente independientes, podemos completarlos hasta una base de $T_pM$:

\[
X_1(p),\ldots,X_k(p),v_{k+1},\ldots,v_n.
\]
Tomamos una carta centrada en $p$,

\[
(V,y=(y^1,\ldots,y^n)),\qquad y(p)=0,
\]
tal que

\[
\left.\frac{\partial}{\partial y^i}\right|_p=X_i(p),\qquad i=1,\ldots,k.
\]
Esto se obtiene realizando un cambio lineal de coordenadas en una carta cualquiera centrada en $p$. 2. Consideramos el subconjunto

\[
S=\{q\in V:y^1(q)=\cdots=y^k(q)=0\}.
\]
Definimos

\[
W=\{(r_{k+1},\ldots,r_n)\in\mathbb R^{n-k}:(0,\ldots,0,r_{k+1},\ldots,r_n)\in y(V)\}.
\]
Entonces la aplicación

\[
\eta:W\longrightarrow S\subseteq M
\]
dada por

\[
\eta(r_{k+1},\ldots,r_n)=y^{-1}(0,\ldots,0,r_{k+1},\ldots,r_n)
\]
es una parametrización local de $S$. En particular,

\[
\eta(0)=p,
\]
porque $y(p)=0$. 3. Para $j\in\{k+1,\ldots,n\}$ probemos que

\[
d\eta_0\left(\left.\frac{\partial}{\partial r_j}\right|_0\right)=\left.\frac{\partial}{\partial y^j}\right|_p.
\]
Sea

\[
v=d\eta_0\left(\left.\frac{\partial}{\partial r_j}\right|_0\right)\in T_pM.
\]
Escribimos $v$ en la base coordenada:

\[
v=\sum_{\ell=1}^n a^\ell\left.\frac{\partial}{\partial y^\ell}\right|_p.
\]
Evaluando ambos lados en la función coordenada $y^m$, obtenemos

\[
v(y^m)=\sum_{\ell=1}^n a^\ell\left.\frac{\partial}{\partial y^\ell}\right|_p(y^m)=\sum_{\ell=1}^n a^\ell\delta_\ell^m=a^m.
\]
Por otro lado, por definición del diferencial,

\[
v(y^m)=d\eta_0\left(\left.\frac{\partial}{\partial r_j}\right|_0\right)(y^m)=\left.\frac{\partial}{\partial r_j}\right|_0(y^m\circ\eta).
\]
Como

\[
y\circ\eta(r_{k+1},\ldots,r_n)=(0,\ldots,0,r_{k+1},\ldots,r_n),
\]
se tiene

\[
y^m\circ\eta=\begin{cases}0,&1\leq m\leq k,\\r_m,&k+1\leq m\leq n.\end{cases}
\]
Por tanto,

\[
a^m=\left.\frac{\partial}{\partial r_j}\right|_0(y^m\circ\eta)=\delta_j^m.
\]
Luego todos los coeficientes son cero salvo $a^j=1$, y así

\[
d\eta_0\left(\left.\frac{\partial}{\partial r_j}\right|_0\right)=\left.\frac{\partial}{\partial y^j}\right|_p.
\]
\begin{enumerate}[start=4]
\item Sea $\theta^i$ el flujo local de $X_i$. Como
\end{enumerate}
\[
[X_i,X_j]=0,
\]
los flujos locales conmutan:

\[
\theta_s^i\circ\theta_t^j=\theta_t^j\circ\theta_s^i
\]
siempre que ambas composiciones estén definidas. Restringiendo los parámetros a un entorno suficientemente pequeño $A\subseteq\mathbb R^n$ de $0$, definimos

\[
\sigma:A\longrightarrow M
\]
por

\[
\sigma(r_1,\ldots,r_n)=\theta_{r_1}^1\circ\cdots\circ\theta_{r_k}^k\bigl(\eta(r_{k+1},\ldots,r_n)\bigr).
\]
Además,

\[
\sigma(0)=p,
\]
porque $\theta_0^i=\operatorname{Id}_M$ para todo $i$. 5. Para $1\leq i\leq k$, definimos la curva

\[
c_i:\mathbb R\longrightarrow\mathbb R^n,\qquad c_i(t)=(0,\ldots,0,t,0,\ldots,0),
\]
donde $t$ ocupa la posición $i$. Su vector velocidad en $0$ es

\[
c_i'(0)=(dc_i)_0\left(\left.\frac{\partial}{\partial t}\right|_0\right).
\]
Probemos que

\[
c_i'(0)=\left.\frac{\partial}{\partial r_i}\right|_0.
\]
Escribimos

\[
c_i'(0)=\sum_{\ell=1}^n b^\ell\left.\frac{\partial}{\partial r_\ell}\right|_0.
\]
Evaluando en la función coordenada $r_m$, obtenemos

\[
b^m=c_i'(0)(r_m).
\]
Por definición de velocidad,

\[
c_i'(0)(r_m)=(dc_i)_0\left(\left.\frac{\partial}{\partial t}\right|_0\right)(r_m)=\left.\frac{\partial}{\partial t}\right|_0(r_m\circ c_i)=\left.\frac{d}{dt}\right|_{t=0}(r_m\circ c_i)(t).
\]
Pero

\[
(r_m\circ c_i)(t)=\begin{cases}t,&m=i,\\0,&m\neq i.\end{cases}
\]
Por tanto,

\[
b^m=\delta_i^m,
\]
y así

\[
c_i'(0)=\left.\frac{\partial}{\partial r_i}\right|_0.
\]
\begin{enumerate}[start=6]
\item Aplicando la propiedad del diferencial respecto de velocidades,
\end{enumerate}
\[
d\sigma_0(c_i'(0))=(\sigma\circ c_i)'(0).
\]
Como

\[
c_i'(0)=\left.\frac{\partial}{\partial r_i}\right|_0,
\]
obtenemos

\[
d\sigma_0\left(\left.\frac{\partial}{\partial r_i}\right|_0\right)=(\sigma\circ c_i)'(0).
\]
Ahora,

\[
\sigma(c_i(t))=\sigma(0,\ldots,t,\ldots,0)=\theta_t^i(p).
\]
Por consiguiente,

\[
d\sigma_0\left(\left.\frac{\partial}{\partial r_i}\right|_0\right)=\left.\frac{d}{dt}\right|_{t=0}\theta_t^i(p).
\]
Como $t\mapsto\theta_t^i(p)$ es la curva integral de $X_i$ que parte de $p$,

\[
\frac{d}{dt}\theta_t^i(p)=X_i(\theta_t^i(p)).
\]
Evaluando en $t=0$,

\[
\left.\frac{d}{dt}\right|_{t=0}\theta_t^i(p)=X_i(\theta_0^i(p))=X_i(p).
\]
Luego

\[
d\sigma_0\left(\left.\frac{\partial}{\partial r_i}\right|_0\right)=X_i(p),\qquad i=1,\ldots,k.
\]
\begin{enumerate}[start=7]
\item Para $j=k+1,\ldots,n$, si solo variamos $r_j$, todos los parámetros de los flujos son cero. Por tanto,
\end{enumerate}
\[
\sigma(0,\ldots,0,r_{k+1},\ldots,r_n)=\eta(r_{k+1},\ldots,r_n).
\]
De aquí,

\[
d\sigma_0\left(\left.\frac{\partial}{\partial r_j}\right|_0\right)=d\eta_0\left(\left.\frac{\partial}{\partial r_j}\right|_0\right)=\left.\frac{\partial}{\partial y^j}\right|_p.
\]
\begin{enumerate}[start=8]
\item Por los pasos anteriores, $d\sigma_0$ transforma la base canónica de $T_0\mathbb R^n$ en
\end{enumerate}
\[
X_1(p),\ldots,X_k(p),\left.\frac{\partial}{\partial y^{k+1}}\right|_p,\ldots,\left.\frac{\partial}{\partial y^n}\right|_p.
\]
Esta es una base de $T_pM$. Por tanto,

\[
d\sigma_0:T_0\mathbb R^n\longrightarrow T_pM
\]
es un isomorfismo. Por el teorema de la función inversa, después de restringir $A$, existen un entorno $A_0\subseteq\mathbb R^n$ de $0$ y un entorno $U\subseteq M$ de $p$ tales que

\[
\sigma:A_0\longrightarrow U
\]
es un difeomorfismo. Definimos la carta

\[
\varphi=\sigma^{-1}:U\longrightarrow A_0,\qquad \varphi=(x^1,\ldots,x^n).
\]
\begin{enumerate}[start=9]
\item Falta probar que, en estas coordenadas,
\end{enumerate}
\[
X_i=\frac{\partial}{\partial x^i},\qquad i=1,\ldots,k.
\]
Sea

\[
r=(r_1,\ldots,r_n)\in A_0
\]
y sea $q=\sigma(r)$. Para $t$ suficientemente pequeño, usando que los flujos conmutan,

\[
\begin{aligned}\sigma(r_1,\ldots,r_i+t,\ldots,r_n)&=\theta_{r_1}^1\circ\cdots\circ\theta_{r_i+t}^i\circ\cdots\circ\theta_{r_k}^k\bigl(\eta(r_{k+1},\ldots,r_n)\bigr)\\&=\theta_t^i\left(\theta_{r_1}^1\circ\cdots\circ\theta_{r_i}^i\circ\cdots\circ\theta_{r_k}^k\bigl(\eta(r_{k+1},\ldots,r_n)\bigr)\right)\\&=\theta_t^i(\sigma(r)).\end{aligned}
\]
Definimos

\[
\gamma(t)=r+te_i.
\]
Entonces

\[
\gamma(0)=r
\]
y

\[
\gamma'(0)=\left.\frac{\partial}{\partial r_i}\right|_r.
\]
Por tanto,

\[
\begin{aligned}d\sigma_r\left(\left.\frac{\partial}{\partial r_i}\right|_r\right)&=d\sigma_r(\gamma'(0))\\&=(\sigma\circ\gamma)'(0)\\&=\left.\frac{d}{dt}\right|_{t=0}\theta_t^i(\sigma(r))\\&=X_i(\sigma(r)).\end{aligned}
\]
Por definición del campo coordenado asociado a la carta $\varphi=\sigma^{-1}$,

\[
\left.\frac{\partial}{\partial x^i}\right|_{\sigma(r)}=d\sigma_r\left(\left.\frac{\partial}{\partial r_i}\right|_r\right).
\]
Luego

\[
\left.\frac{\partial}{\partial x^i}\right|_{\sigma(r)}=X_i(\sigma(r)).
\]
Como todo punto de $U$ es de la forma $\sigma(r)$, concluimos que

\[
X_i|_U=\frac{\partial}{\partial x^i},\qquad i=1,\ldots,k.
\]
Esto completa la demostración.

\bookqed
\end{bookproof}
